\section{Virtual threads, reativo ou worker threads}
\label{sec:vt-vs-reativo}

Por baixo, o Quarkus tem um motor reativo: o guia
\citetitle{doc:quarkus-reactive-architecture} diz que ele é baseado no
Eclipse Vert.x e no Netty e cuida do I/O não bloqueante, com poucas
threads de I/O (\termo{event loop}{Laço de eventos: thread que atende
muitas conexões alternando entre elas sem bloquear; se ela bloqueia,
todas as conexões que ela atende param.}) que atendem muitas conexões
e por isso precisam ser usadas com cuidado. Sobre esse motor, o guia
\citetitle{doc:quarkus-virtual-threads} compara três modelos que podem
conviver na mesma aplicação:

\small
\begin{longtable}{@{}p{3.6cm} p{4.2cm} p{4.4cm}@{}}
\toprule
Modelo & Vantagem & Custo, segundo o guia \\
\midrule \endhead
Código síncrono em \emph{worker thread} & Código simples & A
  concorrência fica limitada pelo número de workers \\
Código reativo no event loop (\texttt{Uni}, \texttt{Multi}) & Alta
  concorrência e pouco uso de recursos & Código mais complexo \\
Código síncrono em virtual thread (\texttt{@RunOn\-Virtual\-Thread}) &
  Código simples & Risco de \emph{pinning}, monopolização e pool de
  objetos ineficiente \\
\bottomrule
\end{longtable}
\normalsize

\paragraph{O que a Parte~II já explicou.}
A seção~\ref{sec:virtual-threads} mostrou que virtual threads servem para
escala em I/O, não para velocidade; que o pinning em
\texttt{synchronized} acabou no JDK~24 (o guia do Quarkus confirma: em
Java~24+, \texttt{synchronized} \enquote{no longer cause pinning}); e
que pool via \texttt{ThreadLocal} é contraproducente. O guia acrescenta
três pontos do ponto de vista do Quarkus:
\begin{itemize}
\item \texttt{@RunOnVirtualThread} só vale em endpoints bloqueantes, por
  \texttt{@Blocking} ou pela assinatura (seção~\ref{sec:panache-rest});
  um método que devolve \texttt{Uni} já roda no event loop.
\item Monopolização: o escalonador atual não preempta tarefas, então uma
  computação longa segura a carrier; o guia avisa que, em ambiente
  restrito, isso pode levar a falta de memória e ao container ser
  encerrado. Para trabalho de CPU, o guia chama virtual threads de
  \enquote{useless and counterproductive}.
\item O ecossistema Java ainda não está todo pronto. Drivers JDBC do
  PostgreSQL anteriores ao 42.6.0 causam pinning frequente em Java~21 a
  23; em Java~24+, a maioria dos drivers JDBC não prende mais a carrier.
  Clientes Vert.x com Mutiny oferecem métodos \texttt{andAwait()} que
  bloqueiam sem prender a carrier.
\end{itemize}

\subsection*{A decisão}

\begin{description}
\item[Quando escolher virtual threads] Serviço de I/O com muitas
  requisições concorrentes que esperam banco ou HTTP, time confortável
  com código imperativo, JDK~24 ou mais novo (ou JDK~21 com as
  bibliotecas conferidas para pinning). É o caminho de menor mudança a
  partir de um serviço bloqueante.
\item[Quando escolher reativo] Streaming, muitas conexões longas
  (WebSocket, SSE), composição de várias chamadas assíncronas, ou quando
  as bibliotecas que você usa já são reativas. Exige que o time saiba não
  bloquear o event loop.
\item[Quando ficar nas worker threads] Concorrência moderada e
  previsível, código legado com \texttt{ThreadLocal} pesado, ou trabalho
  de CPU. O limite de workers funciona como proteção natural.
\item[O que medir] Sob o teste de carga da Parte~I: p99 e throughput
  nos três modelos, número de threads de plataforma e RSS. Com JFR
  (seção~\ref{sec:ferramentas}), o evento \texttt{jdk.VirtualThreadPinned}
  mostra pinning residual. Monitore também as conexões do pool do banco:
  virtual threads não criam conexões, e milhares de requisições
  concorrentes passam a disputar as mesmas poucas.
\end{description}

O repositório \citetitle{doc:quarkus-perf-comparison} tem uma variante
\texttt{quarkus3-virtual} da mesma aplicação, um bom ponto de partida
para comparar os modelos com números do seu hardware.
